\section{Viscous dissipation and material deformation}\label{r:joined-sec:metric}

The deformation of a material parcel is controlled in an integral
sense by viscous dissipation. Let $X(a,t)$ follow the parcel initially labeled $a$, let $F=\nabla_aX$, and let $C=F^TF$. The positive matrix $C$ records the lengths and angles of infinitesimal material vectors. Incompressibility preserves volume, while the symmetric velocity gradient $S=(\nabla u+\nabla u^T)/2$ changes this metric:
\begin{equation}\label{r:joined-eq:metric}
\dot F=(\nabla u)(X,t)F,\qquad
\det F=1,\qquad
\dot C=2F^TS(X,t)F.
\end{equation}
The full metric speed has the exact squared norm
\begin{equation}\label{r:joined-eq:metricspeed}
\tr(C^{-1}\dot C\,C^{-1}\dot C)=4|S(X,t)|^2.
\end{equation}
Indeed, $C^{-1}\dot C=2F^{-1}SF$, and cyclicity of the trace removes the two deformation matrices. After changing variables by the incompressible flow and integrating by parts in space,
\begin{equation}\label{r:joined-eq:metricaction}
\frac\nu2\int_0^t\!\int_{\R^3}
\tr(C^{-1}\dot C\,C^{-1}\dot C)\,da\,ds
=\nu\int_0^t\norm{\nabla u(s)}_2^2ds.
\end{equation}
Here $\norm S_2^2=\frac12\norm{\nabla u}_2^2$. The dissipation integral consequently controls the rate of change
of the material metric, including shape changes that preserve volume.

The energy budget makes this quantitative through the full time interval. Define
\begin{equation*}
F_T=\int_0^T\norm f_2dt,\quad
M_T=\norm{u_0}_2+F_T,\quad
K_T=\norm{u_0}_2^2+2M_TF_T.
\end{equation*}
The energy identity gives $\sup_t\norm u_2\le M_T$ and
$2\nu\int_0^T\norm{\nabla u}_2^2\le K_T$.
If $\sigma_i$ are the singular values of $F$, then the derivatives of $\log\sigma_i$ are diagonal entries of $S$ in a suitable orthonormal frame. Their squared sum is bounded by $|S|^2$. Integrating from $F(a,0)=I$ gives
\begin{equation*}
\norm{\log C(a,t)}_F\le2\int_0^t|S(X(a,s),s)|\,ds.
\end{equation*}
Repeated singular values are handled by diagonalizing the symmetric variation within each eigenspace. Cauchy--Schwarz and volume preservation now yield the maximal estimate
\begin{equation}\label{r:joined-eq:maxdistortion}
\int_{\R^3}\sup_{t<T}\norm{\log C(a,t)}_F^2\,da
\le\frac{TK_T}{\nu}.
\end{equation}
This is a bound over original material labels with the time supremum inside the integral. It rules out infinite distortion on a set of labels of positive measure. It permits concentration onto label sets whose measure tends to zero. A uniform bound over all labels would require further control beyond \eqref{r:joined-eq:maxdistortion}.

